% =====================================================================
% Problem statement -- carries over the content of the approved topic
% proposal. Please do not change the structure.
% =====================================================================
\vorspannkapitel{Problemstellung}
\abschnittsautor{\alleautoren}

\vorspannabschnitt{Ausgangslage}

Automatisierte Spielanalyse (Laufdistanzen, Positionsdaten, Ereigniserkennung) ist im Profifußball Standard, 
für Amateur- und Jugendvereine aber kaum zugänglich: bestehende 
Systeme sind teuer, an Spezialhardware gebunden und geschlossen. 
Vereine filmen zwar oft mit einfachen Kameras, können die 
Aufnahmen aber nicht automatisiert auswerten.

\vorspannabschnitt{Untersuchungsanliegen der individuellen Themenstellungen}

\vorspannunterpunkt{1) Teilbereich A Akpinar }
Welche Messabweichungen und Genauigkeitsgrenzen entstehen bei der 
2D-Homographie ohne echte 3D-Tiefsensoren

\vorspannunterpunkt{2) Teilbereich B Ishak}
Wie verändert sich die Detektionsgenauigkeit bei Spielerverdeckung
durch spezifisches Fine-Tuning im vergleich zu Standard-Basismodellen?

\vorspannunterpunkt{3) Teilbereich C Rahal}
Wie muss eine Systemarchitektur ausgelegt sein, 
um hochfrequente Tracking-Rohdaten mit minimaler Latenz 
abzufragen und für Laien intuitiv darzustellen?
\vorspannabschnitt{Zielsetzung}

Am Ende der Arbeit liegt ein funktionsfähiger Software-Prototyp 
vor, der Fußballspiele über ein Zwei-Kamera-Panoramasystem 
erfasst, Spieler mittels Computer Vision trackt und deren 
Bewegungsdaten als reale 2D-Spielfeldkoordinaten in einem 
interaktiven Web-Dashboard visualisiert.

\vorspannabschnitt{Geplantes Ergebnis der individuellen Themenstellungen}

\vorspannunterpunkt{1) Teilbereich A Akpinar}
Artefakt:

Eine Software-Pipeline, die das Video-Stitching beider Kameras 
durchführt, Tracking-IDs vergibt und die Pixelkoordinaten über 
eine kalibrierte Homographie-Matrix in ein reales Koordinatensystem
(in Metern) umrechnet.

Messbarkeit:

Toleranz- und Fehleranalyse: Abweichung der berechneten Distanzen
gegenüber real vermessenen Referenzstrecken auf dem Platz 
(z. B. prozentualer Messfehler bei 16-Meter-Raum, Mittellinie) 
sowie Stabilität des Trackings (Anzahl der ID-Switches).

\vorspannunterpunkt{2) Teilbereich B Ishak}
Artefakt:

Ein exportiertes, lauffähiges Modell 
(.pt / .yaml), die zugehörige Trainings- und Inferenz-Pipeline 
sowie der vorbereitete, annotierte Validierungsdatensatz. 

Messbarkeit: 

Quantitative Evaluierungsmetriken: $mAP_{50}$, $mAP_{50-95}$, 
Precision, Recall und Inferenzzeit (FPS / Latenz in ms) bei 
Szenarien mit Verdeckung im Vergleich zum ungetunten Basismodell.

\vorspannunterpunkt{3) Teilbereich C Rahal}
Artefakt:

Ein lauffähiges Softwaresystem bestehend aus 
relationalem/zeitreihenbasiertem Datenbankschema, REST- bzw. 
WebSocket-API zur Datenabfrage und einem interaktiven Web-Frontend.
Wie wird es überprüft 

Messbarkeit:

Funktionale Abnahmekriterien: End-to-End-Test 
Import von Trackingdaten $\rightarrow$ korrekte 
visuelle Ausgabe auf dem 2D-Taktikfeld, Rendern von Heatmaps) 
und Performance-Benchmarks (z. B. API-Antwortzeiten unter Last).

\vorspannabschnitt{Meilensteine}

% Deliberately not a table float: the milestones belong directly under the
% heading, not somewhere else in the document.
\begin{tabular}{@{}p{0.68\linewidth}p{0.24\linewidth}@{}}
    \toprule
    \bfseries Meilenstein & \bfseries Datum \\
    \midrule
    Analyse der Ausgangslage und Anforderungen & 30.09.2026 \\
    Erstellung der Konzepte                    & 31.10.2026 \\
    Implementierung der Systemkomponenten      & 30.11.2026 \\
    Test der Systemkomponenten                 & 15.01.2027 \\
    Integration des Gesamtsystems              & 15.02.2027 \\
    Verfassen der Dokumentation                & 15.03.2027 \\
  Abgabe der Diplomarbeit                    & 01.04.2027 \\
  \bottomrule
\end{tabular}
